% GENERATED FILE. Edit canonical lessons, then run npm run generate:books.
% canonical-source-hash: fnv1a64:3efc1ee67d1cbc7f
% canonical-lessons: IT-C439-orientarsi, IT-C439-smarrirsi, IT-C439-cartina, IT-C439-piantina, IT-C439-bussola

\chapter{Orientarsi}
\label{ch:it-orientarsi}
\hlchaptermodality{439}

\begin{chapteropening}
\textbf{By the end of this chapter:} \emph{I can find my bearings: finding my way, getting lost, the map, the street plan and the compass.}

Complete the last lesson of chapter 439: I can find my bearings: finding my way, getting lost, the map, the street plan and the compass.
\end{chapteropening}

\section[orientarsi]{orientarsi — to find one's way, to get one's bearings}

\label{lesson:IT-C439-orientarsi}



\begin{quote}
Before the new one: say the Italian for to put right, then the Italian for to compensate.
\end{quote}

\subsection*{You'll want to know: orientarsi}
\textbf{orientarsi} — \textquotedblleft{}to find one's way, to get one's bearings\textquotedblright{}.

\begin{quote}
In questa città è facile orientarsi.
\end{quote}

\begin{quote}
\emph{In this city it is easy to find your way.}
\end{quote}

Built on \emph{oriente}, the east: once you knew where the sun rose, you knew the rest.

\begin{cousinweb}
\emph{orientarsi} is built on \emph{oriente} 'east', which is inherited from Latin \emph{oriēns} 'rising (sun)' (EtymDB), from \emph{orīrī} 'to rise'.
\end{cousinweb}

\begin{grammarlens}[title={What you've built}]
One more word for telling, explaining and describing at B1.
\end{grammarlens}

\subsection*{Your turn}
\begin{itemize}
\raggedright
  \item \emph{Say it:} \emph{orientarsi}
  \item \emph{Say it:} \emph{orientarsi}, once more
  \item \emph{Say it:} a sentence of your own with \emph{orientarsi}
  \item \emph{Recall:} say the Italian for to put right, then the Italian for to compensate, then say \emph{orientarsi} again
\end{itemize}

\begin{tcolorbox}[breakable,colback=teal!4,colframe=teal!35!black,title={Before you move on}]
What does orientarsi mean? (\textquotedblleft{}To find one's way, to get one's bearings\textquotedblright{}.) Say it once more.
\end{tcolorbox}
\section[smarrirsi]{smarrirsi — to get lost, to go astray}

\label{lesson:IT-C439-smarrirsi}



\begin{quote}
Before the new one: say the Italian for to compensate, then the Italian for to find one's way.
\end{quote}

\subsection*{You'll want to know: smarrirsi}
\textbf{smarrirsi} — \textquotedblleft{}to get lost, to go astray\textquotedblright{}.

\begin{quote}
Senza la mappa ci siamo smarriti.
\end{quote}

\begin{quote}
\emph{Without the map we got lost.}
\end{quote}

A little more literary than \emph{perdersi}; \emph{oggetti smarriti} — lost property.

\begin{cousinweb}
\emph{smarrirsi} was borrowed from Old French \emph{esmarrir} 'to trouble, to dismay' (EtymDB), which came from a Germanic word.
\end{cousinweb}

\begin{grammarlens}[title={What you've built}]
One more word for telling, explaining and describing at B1.
\end{grammarlens}

\subsection*{Your turn}
\begin{itemize}
\raggedright
  \item \emph{Say it:} \emph{smarrirsi}
  \item \emph{Say it:} \emph{smarrirsi}, once more
  \item \emph{Say it:} a sentence of your own with \emph{smarrirsi}
  \item \emph{Recall:} say the Italian for to compensate, then the Italian for to find one's way, then say \emph{smarrirsi} again
\end{itemize}

\begin{tcolorbox}[breakable,colback=teal!4,colframe=teal!35!black,title={Before you move on}]
What does smarrirsi mean? (\textquotedblleft{}To get lost, to go astray\textquotedblright{}.) Say it once more.
\end{tcolorbox}
\section[la cartina]{la cartina — the map (of a town or an area)}

\label{lesson:IT-C439-cartina}



\begin{quote}
Before the new one: say the Italian for to find one's way, then the Italian for to get lost.
\end{quote}

\subsection*{You'll want to know: la cartina}
\textbf{la cartina} — \textquotedblleft{}the map (of a town or an area)\textquotedblright{}.

\begin{quote}
All'ufficio turistico danno una cartina gratis.
\end{quote}

\begin{quote}
\emph{At the tourist office they give you a free map.}
\end{quote}

A little \emph{carta}; \emph{la cartina stradale} — road map.

\begin{cousinweb}
\emph{cartina} is a little \emph{carta} 'paper, map', which is inherited from Latin \emph{charta}, from Greek \emph{khártēs} 'papyrus' (EtymDB).
\end{cousinweb}

\begin{grammarlens}[title={What you've built}]
One more word for telling, explaining and describing at B1.
\end{grammarlens}

\subsection*{Your turn}
\begin{itemize}
\raggedright
  \item \emph{Say it:} \emph{la cartina}
  \item \emph{Say it:} \emph{la cartina}, once more
  \item \emph{Say it:} a sentence of your own with \emph{la cartina}
  \item \emph{Recall:} say the Italian for to find one's way, then the Italian for to get lost, then say \emph{la cartina} again
\end{itemize}

\begin{tcolorbox}[breakable,colback=teal!4,colframe=teal!35!black,title={Before you move on}]
What does la cartina mean? (\textquotedblleft{}The map (of a town or an area)\textquotedblright{}.) Say it once more.
\end{tcolorbox}
\section[la piantina]{la piantina — the street plan; the floor plan}

\label{lesson:IT-C439-piantina}



\begin{quote}
Before the new one: say the Italian for to get lost, then the Italian for the map.
\end{quote}

\subsection*{You'll want to know: la piantina}
\textbf{la piantina} — \textquotedblleft{}the street plan; the floor plan\textquotedblright{}.

\begin{quote}
Sulla piantina il museo è vicino al fiume.
\end{quote}

\begin{quote}
\emph{On the plan the museum is near the river.}
\end{quote}

Also a little plant: \emph{una piantina di basilico}.

\begin{cousinweb}
\emph{piantina} is a little \emph{pianta} 'plan, plant', which is inherited from Latin \emph{planta} 'sole of the foot' (EtymDB): a plan is the footprint of a place.
\end{cousinweb}

\begin{grammarlens}[title={What you've built}]
One more word for telling, explaining and describing at B1.
\end{grammarlens}

\subsection*{Your turn}
\begin{itemize}
\raggedright
  \item \emph{Say it:} \emph{la piantina}
  \item \emph{Say it:} \emph{la piantina}, once more
  \item \emph{Say it:} a sentence of your own with \emph{la piantina}
  \item \emph{Recall:} say the Italian for to get lost, then the Italian for the map, then say \emph{la piantina} again
\end{itemize}

\begin{tcolorbox}[breakable,colback=teal!4,colframe=teal!35!black,title={Before you move on}]
What does la piantina mean? (\textquotedblleft{}The street plan; the floor plan\textquotedblright{}.) Say it once more.
\end{tcolorbox}
\section[la bussola]{la bussola — the compass}

\label{lesson:IT-C439-bussola}



\begin{quote}
Before the new one: say the Italian for the map, then the Italian for the street plan.
\end{quote}

\subsection*{You'll want to know: la bussola}
\textbf{la bussola} — \textquotedblleft{}the compass\textquotedblright{}.

\begin{quote}
In montagna porto la bussola.
\end{quote}

\begin{quote}
\emph{In the mountains I carry a compass.}
\end{quote}

\emph{Perdere la bussola} — to lose your head.

\begin{cousinweb}
\emph{bussola} is inherited from Late Latin \emph{buxula} 'little box' (EtymDB), from \emph{buxus} 'boxwood': the needle sat in a small wooden box.
\end{cousinweb}

\begin{grammarlens}[title={What you've built}]
One more word for telling, explaining and describing at B1.
\end{grammarlens}

\subsection*{Your turn}
\begin{itemize}
\raggedright
  \item \emph{Say it:} \emph{la bussola}
  \item \emph{Say it:} \emph{la bussola}, once more
  \item \emph{Say it:} a sentence of your own with \emph{la bussola}
  \item \emph{Recall:} say the Italian for the map, then the Italian for the street plan, then say \emph{la bussola} again
\end{itemize}

\begin{tcolorbox}[breakable,colback=teal!4,colframe=teal!35!black,title={Before you move on}]
What does la bussola mean? (\textquotedblleft{}The compass\textquotedblright{}.) Say it once more.
\end{tcolorbox}
